\documentclass[11pt]{article}
\usepackage[T1]{fontenc}
\usepackage{lmodern}
\usepackage[margin=1in]{geometry}
\usepackage{amsmath,amssymb,amsthm,mathtools}
\usepackage{enumitem}
\usepackage{microtype}
\usepackage{xcolor}
\usepackage[colorlinks=true,linkcolor=blue!45!black,citecolor=blue!45!black,urlcolor=blue!45!black]{hyperref}
\hypersetup{pdftitle={Critical-boundary trace limits for zeta-phase Sonin operators},
pdfauthor={Drafted for Edward Baker},
pdfsubject={Sonin projections, zeta phases, uniform trace bounds, and projection loss}}
\numberwithin{equation}{section}
\newtheorem{theorem}{Theorem}[section]
\newtheorem{proposition}[theorem]{Proposition}
\newtheorem{lemma}[theorem]{Lemma}
\newtheorem{corollary}[theorem]{Corollary}
\theoremstyle{definition}
\newtheorem{definition}[theorem]{Definition}
\theoremstyle{remark}
\newtheorem{remark}[theorem]{Remark}
\newcommand{\HH}{\mathcal H}
\newcommand{\Sch}{\mathcal S}
\newcommand{\Tr}{\operatorname{Tr}}
\newcommand{\Ran}{\operatorname{Ran}}
\newcommand{\Rea}{\operatorname{Re}}
\newcommand{\Ima}{\operatorname{Im}}
\newcommand{\ii}{\mathrm i}
\newcommand{\Id}{I}
\newcommand{\J}{\mathcal J}
\newcommand{\Fcal}{\mathcal F}
\newcommand{\Qcal}{\mathcal Q}
\newcommand{\Sone}{\mathfrak S_1}
\newcommand{\Stwo}{\mathfrak S_2}
\newcommand{\norm}[1]{\left\lVert #1\right\rVert}
\newcommand{\abs}[1]{\left\lvert #1\right\rvert}
\newcommand{\inner}[2]{\left\langle #1,#2\right\rangle}
\newcommand{\dd}{\,\mathrm d}
\newcommand{\mucrit}{\mu_{\mathrm{crit}}}
\newcommand{\Zcrit}{Z_{\mathrm{crit}}}
\newcommand{\Zoff}{Z_{\mathrm{off}}}
\setlist[enumerate]{label=\textup{(\roman*)},itemsep=3pt}
\setlength{\emergencystretch}{2em}
\title{Critical-boundary trace limits\\for zeta-phase Sonin operators}
\author{Drafted for Edward Baker}
\date{30 September 2026\\[3pt]\small Working draft}
\begin{document}
\maketitle

\begin{abstract}
We study the Sonin-type intersection projections associated with the
unimodular phases
\(v_\sigma(t)=e^{2\ii\theta(t)}
\zeta(\sigma+\ii t)/\zeta(\sigma-\ii t)\), \(\sigma>1/2\).
A boundary-resolvent formula separates the positive projection trace from
the phase derivative. We prove uniform logarithmic-square bounds for the
unit-window masses of the projection, a positive majorant, and the
compressed cutoff square. These bounds give finite traces and uniform
frequency tails for every compact smooth source. As \(\sigma\downarrow1/2\),
the majorant and every fixed Abel regularization converge, after source
smoothing, to the counting measure of critical-line zeros. Every
subsequential limit of the exact projection measure is supported on those
zeros, with weights between zero and their multiplicities. The remaining
projection defect is positive and concentrates at cutoff spectral value
one. An explicit rational phase has the same local concentration and a
positive gap at every prelimit parameter, but its exact intersection
projection vanishes identically. Finally, a contour comparison retains
the pole and all off-line zero crossings, distinguishing the critical-zero
measure from the complete Weil form. Neither survival of all critical
mass in the exact zeta projection nor the Riemann hypothesis is proved.
\end{abstract}

\noindent\textbf{Draft status.}
This manuscript consolidates internal research derivations.
Independent specialist verification and a full priority assessment remain
outstanding. All results below are analytic; no numerical certificate is
used in their proofs.

\tableofcontents

\section{Introduction and main results}\label{sec:intro}

A positive trace attached to a Hilbert-space compression is a natural
candidate for comparison with Weil's quadratic form. The difficult part
of such a comparison is the identification of the arithmetic functional.
An orthogonal projection may converge strongly to zero while its
smoothed trace retains a nonzero limit. Conversely, concentration in a
phase derivative need not survive passage to an exact intersection
kernel. This paper separates these phenomena for a family of zeta phases.

Connes and Consani express the archimedean Weil distribution as a Sonin
trace with a prolate correction \cite{CCWeil}. Their quasi-inner framework
constructs semilocal Sonin spaces and maps between them \cite{CCQuasi}.
Connes, Consani, and Moscovici further analyze semilocal transport and
the associated Hilbert spaces of entire functions \cite{CCM}.
The relation of Sonine spaces to zeta zeros has earlier developments
in de Branges's framework and Burnol's work \cite{Burnol}. Spectral
interpretations that distinguish critical zeros from possible noncritical
zeros also predate the present construction \cite{ConnesTrace}.

The contribution considered here is the critical-boundary behavior of
one specified phase family: uniform source-trace estimates, a precise
localization of possible projection loss, and a solvable example showing
the limits of local phase information. We use standard two-projection
algebra \cite{Halmos}, Schwartz-kernel smoothing, local estimates for
\(\zeta'/\zeta\), and elementary contour arguments. We do not claim a
new general positivity principle or a new version of the explicit formula.

Here is the operator setting. On \(\HH=L^2(\mathbb R,\dd x)\), let
\[
P=1_{(0,\infty)},\qquad \chi=\Id-P,\qquad
(\J h)(x)=h(-x).
\]
With the Fourier convention fixed in Section~\ref{sec:setup}, set
\begin{equation}\label{eq:phase-intro}
v_\sigma(t)=
\frac{\Gamma(1/4+\ii t/2)}{\Gamma(1/4-\ii t/2)}
\pi^{-\ii t}
\frac{\zeta(\sigma+\ii t)}{\zeta(\sigma-\ii t)},
\qquad
\Fcal_\sigma=\J v_\sigma(D).
\end{equation}
The operator \(\Fcal_\sigma\) is a selfadjoint unitary involution. Define
\begin{equation}\label{eq:operators-intro}
T_\sigma=\chi\Fcal_\sigma P,\quad
C_\sigma=\chi\Fcal_\sigma\chi,\quad
\Pi_\sigma=\operatorname{proj}\ker(T_\sigma|_{P\HH}),
\quad L_\sigma=P-T_\sigma^*T_\sigma.
\end{equation}
All compressed operators are extended by zero on the complementary
half-line. In particular \(0\le\Pi_\sigma\le L_\sigma\le P\).
We call
\[
H_\sigma=L_\sigma-\Pi_\sigma
\]
the \emph{return operator}; this is an independently defined positive
operator.

For a bounded positive operator \(G\), its frequency measure is initially
defined with infinite values allowed by
\begin{equation}\label{eq:measure-intro}
\nu_G(E)=\norm{1_E(D)G^{1/2}}_{\Stwo}^2.
\end{equation}
Local finiteness is part of the theorem. Write
\(\nu_\sigma=\nu_{\Pi_\sigma}\), and let
\[
\mucrit=\sum_{\rho=1/2+\ii\gamma}m_\rho\delta_\gamma
\]
sum over distinct nontrivial critical-line zeros, including both signs
of the ordinates, with multiplicity \(m_\rho\).

\begin{theorem}[Uniform bounds and the critical limit]\label{thm:main}
For every finite \(S>1/2\), there is \(M_S<\infty\) such that
\begin{equation}\label{eq:main-window}
\nu_G([j,j+1])\le M_S\log^2(2+|j|),\qquad j\in\mathbb Z,
\end{equation}
uniformly for \(1/2<\sigma\le S\) and
\(G=\Pi_\sigma,L_\sigma,C_\sigma^2\).
For every \(f\in\Sch(\mathbb R)\), all corresponding positive sandwich
traces are finite, with uniform tails in frequency. Moreover:
\begin{enumerate}
\item The projections \(\Pi_\sigma\) tend strongly to zero as
\(\sigma\downarrow1/2\), but
\[
\Tr(f(D)L_\sigma f(D)^*)\longrightarrow
\int |f|^2\dd\mucrit.
\]
\item With \(R_\sigma[f]=\Tr(f(D)H_\sigma f(D)^*)\ge0\),
\begin{equation}\label{eq:main-defect}
\Tr(f(D)\Pi_\sigma f(D)^*)
=\int|f|^2\dd\mucrit-R_\sigma[f]+o_f(1).
\end{equation}
Every sequence \(\sigma_n\downarrow1/2\) has a subsequence for which
\(\nu_{\sigma_n}\) converges vaguely. Every such limit is
\[
\nu_*=\sum_{\rho=1/2+\ii\gamma}c_\gamma\delta_\gamma,
\qquad 0\le c_\gamma\le m_\rho.
\]
\item For each fixed \(0\le r<1\), the positive contractions
\[
\Pi_{\sigma,r}=P-T_\sigma^*(\Id_\chi-rC_\sigma^2)^{-1}T_\sigma
\]
satisfy
\[
\Tr(f(D)\Pi_{\sigma,r}f(D)^*)\longrightarrow
\int |f|^2\dd\mucrit.
\]
\end{enumerate}
\end{theorem}

The theorem does not assert convergence of the exact projection measures
or equality \(c_\gamma=m_\rho\).
Section~\ref{sec:return} identifies the remaining defect as mass
approaching the endpoint of the spectrum of \(C_\sigma^2\).
Section~\ref{sec:model} gives an explicit family for which every exact
projection is zero although the limiting majorant measure is a unit atom.
It also computes the dependence on a jointly varying Abel parameter.

To state the arithmetic consequence, let \(F\in C_c^\infty(\mathbb R)\),
let \(C_F\) be convolution by \(F\), and put
\[
a_F(z)=\widehat F(z)\overline{\widehat F(\overline z)}.
\]
Write \(B_\sigma[F]=\Tr(C_F\Pi_\sigma C_F^*)\) and
\(\Zcrit[F]=\int|\widehat F|^2\dd\mucrit\).
For the prepared sources
\begin{equation}\label{eq:prepared}
F=(-\partial_x^2+1/4)h,\qquad h\in C_c^\infty(\mathbb R),
\end{equation}
the pole terms in the critical explicit formula vanish.
With \(Q\) normalized in Section~\ref{sec:arithmetic}, the comparison is
\begin{equation}\label{eq:intro-arithmetic}
B_\sigma[F]-Q[F]
=-R_\sigma[\widehat F]-\Zoff[F]+o_F(1),\qquad
\Zoff[F]=2\Rea\sum_{\beta>1/2}m_{\beta+\ii\gamma}
a_F\bigl(\gamma+\ii(\beta-1/2)\bigr).
\end{equation}
At nonreal arguments \(a_F\) is not a modulus square.
Consequently \(\Zoff\) has no asserted sign. The positive limits in
Theorem~\ref{thm:main} count critical-line zeros; unconditional
identification with the complete Weil form is a separate problem.

\section{Conventions, cutoff geometry, and smoothing}\label{sec:setup}

We use inner products antilinear in the first variable and
\begin{equation}\label{eq:fourier}
\widehat h(t)=\int_{\mathbb R}h(x)e^{-\ii tx}\dd x,\qquad
h(x)=\int_{\mathbb R}\widehat h(t)e^{\ii tx}\frac{\dd t}{2\pi}.
\end{equation}
The notation \(b(D)\) denotes the Fourier multiplier with symbol \(b\).
The Riemann--Siegel phase in the abstract is
\[
\theta(t)=\Ima\log\Gamma(1/4+\ii t/2)-(t/2)\log\pi,
\]
with the continuous logarithm normalized at \(t=0\).
The reflection \(\J\) is linear. No complex conjugation is included
in \(\J\). The symbols \(\norm{\cdot}_{\Sone}\) and
\(\norm{\cdot}_{\Stwo}\) denote trace and Hilbert--Schmidt norms.

At a zero or pole on the line \(\Rea s=\sigma\), the numerator and
denominator of the zeta ratio in \eqref{eq:phase-intro} have the same
order as functions of real \(t\). The ratio therefore has a smooth
removable extension near that frequency. We always use this extension,
not a one-sided distributional limit in \(\sigma\).
On the real line
\[
|v_\sigma|=1,\qquad v_\sigma(-t)=\overline{v_\sigma(t)}.
\]
These identities imply that \(\Fcal_\sigma\) is a selfadjoint
unitary involution and preserves real functions.

For part of the argument it is useful to work with an arbitrary smooth
unimodular phase \(v\) having this symmetry. Set \(V=v(D)\),
\(\Fcal=\J V\), and use \(C,T,L,\Pi,H\) without subscripts.
The elementary block identities are
\begin{align}
TT^*&=\Id_\chi-C^2,\label{eq:tt}\\
L&=P V^*\chi V P=(P\Fcal P)^2,\label{eq:Lfactor}\\
C^2&=\chi V^*PV\chi,\label{eq:Cfactor}\\
\Qcal:=V^*PV-P
&=\begin{pmatrix}-L&T^*C\\ CT&C^2\end{pmatrix}
\quad\hbox{on }P\HH\oplus\chi\HH.\label{eq:Qblocks}
\end{align}
Indeed \(V=\J\Fcal\), \(\J P\J=\chi\), and \(\Fcal^2=\Id\).
On \(P\HH\), \(L=\Id_{P\HH}-T^*T\); hence \(0\le L\le P\).
The kernel projection of \(T\) is the spectral projection of \(L\)
at \(1\), so \(L\Pi=\Pi\) and \(0\le\Pi\le L\).

\begin{lemma}[Schwartz cutoff blocks]\label{lem:hankel}
If \(q\in\Sch(\mathbb R)\), then \(Pq(D)\chi\),
\(\chi q(D)P\), and \([P,q(D)]\) are trace class.
Their trace norms depend continuously on Schwartz seminorms of \(q\).
Furthermore,
\begin{equation}\label{eq:comm-energy}
\norm{[P,q(D)]}_{\Stwo}^2
=\frac1{4\pi^2}\iint_{\mathbb R^2}
\frac{|q(t)-q(s)|^2}{(t-s)^2}\dd t\dd s.
\end{equation}
The last identity also applies whenever the right side is finite and
\(q\) is bounded.
\end{lemma}
\begin{proof}
After reflection of the negative coordinate, the kernel of
\(Pq(D)\chi\) is \(\check q(x+y)\) for \(x,y>0\).
Choose a smooth function equal to one on \([0,\infty)\) and zero
on \((-\infty,-1]\) in each variable. Multiplying
\(\check q(x+y)\) by these cutoffs gives a Schwartz kernel on
\(\mathbb R^2\). Expansion in Hermite functions, or repeated application
of the harmonic oscillator in both variables, gives rapidly decaying
matrix coefficients. Thus the extended operator is trace class, with
trace norm controlled by finitely many Schwartz seminorms.
Compression proves the first assertion. This is also the smoothing
argument used in \cite[Appendix D]{CCWeil}.

In Fourier coordinates with measure \(\dd t/(2\pi)\), the kernel of
\(P\) is
\[
\pi\delta(t-s)-\ii\operatorname{pv}\frac1{t-s}.
\]
The commutator cancels the delta term and has kernel
\(-\ii(q(s)-q(t))/(t-s)\). Its squared Hilbert--Schmidt norm gives
\eqref{eq:comm-energy}. Approximation, or the same kernel calculation
in distributions, proves the finite-energy version.
\end{proof}

\begin{lemma}[Localized phase trace]\label{lem:phase-trace}
For \(b\in C_c^\infty(\mathbb R)\), let \(\phi'\) be the smooth
local phase derivative of \(v\), so \(v'=\ii\phi'v\).
The positive operators \(b(D)Lb(D)^*\) and \(b(D)C^2b(D)^*\)
are trace class, as is \(b(D)\Qcal b(D)^*\), and
\begin{equation}\label{eq:bulk-trace}
-\Tr(b(D)\Qcal b(D)^*)=\int |b(t)|^2\phi'(t)\frac{\dd t}{2\pi}.
\end{equation}
Writing \(X_a=Pa(D)\chi\), \(a=|b|^2\), gives
\begin{equation}\label{eq:L-trace}
\Tr(b(D)Lb(D)^*)=
\int a(t)\phi'(t)\frac{\dd t}{2\pi}
+\Tr(b(D)C^2b(D)^*)+2\Rea\Tr_{\chi\HH}(CTX_a).
\end{equation}
\end{lemma}
\begin{proof}
Whenever \(q,qv\) are Schwartz,
\begin{equation}\label{eq:phase-comm}
q(D)\Qcal
=V^*\bigl([P,(qv)(D)]-[P,q(D)]V\bigr)\in\Sone.
\end{equation}
Apply this first with \(q=b\). Commuting \(b(D)\) past \(P\)
and \(\chi\), using Lemma~\ref{lem:hankel}, proves that \(b(D)\)
times every block in \eqref{eq:Qblocks} is trace class. In particular
the two positive sandwiches are trace class.

The Fourier kernel of \(\Qcal\) is
\[
-\ii\frac{\overline{v(t)}v(s)-1}{t-s}.
\]
Its diagonal value is \(\ii\overline v(t)v'(t)=-\phi'(t)\).
The kernel of its sandwich by \(b\) is smooth and compactly supported
in both frequency variables. The trace is therefore its diagonal
integral, proving \eqref{eq:bulk-trace}.
Taking traces of \eqref{eq:Qblocks} yields \eqref{eq:L-trace}.
The crossing term can be cycled to the negative half-line because
\(X_a\in\Sone\) and all other factors are bounded.
\end{proof}

\begin{remark}
For \(G\ge0\), an orthonormal-basis expansion of
\eqref{eq:measure-intro} and Tonelli's theorem give countable additivity
and
\[
\Tr(b(D)Gb(D)^*)=\int|b|^2\dd\nu_G.
\]
At a parameter where the measure is locally finite it is absolutely
continuous: a Lebesgue-null frequency set has zero spectral projection
\(1_E(D)\). This does not prevent an atomic limit of the measures.
\end{remark}

\section{Actual projections and boundary formulas}\label{sec:boundary}

\subsection{The bounded-transport region}

Let \(\Fcal_\infty=\J v_\infty(D)\), where
\[
v_\infty(t)=
\pi^{-\ii t}\frac{\Gamma(1/4+\ii t/2)}
{\Gamma(1/4-\ii t/2)}.
\]
This is the logarithmic form of the unitary cosine transform at physical
cutoff one. Its compressed negative-half-line operator
\(C_\infty=\chi\Fcal_\infty\chi\) is compact and
\begin{equation}\label{eq:c0}
c_0:=\norm{C_\infty}<1.
\end{equation}
One way to see strictness is to return to the finite physical interval:
the cosine kernel is smooth there and hence compact. Equality of its
norm to one would give a nonzero function supported in that interval
whose cosine transform is also supported there. The transform of a
compactly supported function is entire, and vanishing on the exterior
interval forces it to be zero. Compactness makes this exclusion of an
extremal vector sufficient for \eqref{eq:c0}; compare \cite{CCWeil}.

For \(\sigma>1\), put \(U_ah(x)=h(x-a)\) and
\begin{equation}\label{eq:transport}
D_\sigma=\sum_{n\ge1}\mu(n)n^{-\sigma}U_{\log n},
\qquad
D_\sigma^{-1}=\sum_{n\ge1}n^{-\sigma}U_{\log n}.
\end{equation}
Both series converge absolutely in operator norm; their multipliers
are \(1/\zeta(\sigma+\ii t)\) and \(\zeta(\sigma+\ii t)\).
Both preserve \(P\HH\). Set
\[
\ell_\sigma=\zeta(\sigma)^{-1},\qquad
u_\sigma=\frac{\zeta(\sigma)}{\zeta(2\sigma)},\qquad
D_+=D_\sigma|_{P\HH},\quad
D_-=\chi D_\sigma\chi|_{\chi\HH}.
\]
Then \(\norm{D_+}\le u_\sigma\) and
\(\norm{D_-^{-1}}\le\ell_\sigma^{-1}\).
The inverse assertion for \(D_-\) follows by compressing
\(D_\sigma^{-1}\): the triangularity
\(\chi D_\sigma P=\chi D_\sigma^{-1}P=0\) makes the two products identities.

\begin{proposition}[Gap and exact metric]\label{prop:gap}
For \(\sigma>1\),
\begin{equation}\label{eq:gap}
T_\sigma=D_-T_\infty D_+^{-1},\qquad
\Id_\chi-C_\sigma^2\ge
g_\sigma\Id_\chi,\quad
g_\sigma=(\ell_\sigma/u_\sigma)^2(1-c_0^2)>0,
\end{equation}
and
\begin{equation}\label{eq:projection}
\Pi_\sigma=P-T_\sigma^*(\Id_\chi-C_\sigma^2)^{-1}T_\sigma.
\end{equation}
It is the ordinary orthogonal projection onto
\(D_\sigma\Ran\Pi_\infty\).
\end{proposition}
\begin{proof}
The multiplier relation gives
\(\Fcal_\sigma=D_\sigma\Fcal_\infty D_\sigma^{-1}\).
Compress and use the triangularity just noted to obtain
the formula for \(T_\sigma\).
For \(y\in\chi\HH\),
\[
\norm{T_\sigma^*y}
\ge u_\sigma^{-1}\norm{T_\infty^*D_-^*y}
\ge u_\sigma^{-1}(1-c_0^2)^{1/2}\ell_\sigma\norm y.
\]
This proves the gap. Thus \(T_\sigma\) is onto, and the usual kernel
projection formula gives \eqref{eq:projection}.
The factorization of \(T_\sigma\) identifies its kernel.

Equivalently, if
\(A_\sigma=\Pi_\infty D_\sigma^*D_\sigma\Pi_\infty\)
on \(\Ran\Pi_\infty\), then
\[
\Pi_\sigma=D_\sigma\Pi_\infty A_\sigma^{-1}
\Pi_\infty D_\sigma^*.
\]
The inverse compressed metric is necessary in this formula.
\end{proof}

\subsection{A boundary resolvent with a finite crossing strip}

For \(F\in C_c^\infty\), put \(a=|\widehat F|^2\), and
\[
a_e(t)=\frac{a(t)+a(-t)}2,\qquad X_{F,e}=Pa_e(D)\chi.
\]
The phase derivative is
\begin{equation}\label{eq:phase-deriv}
\phi_\sigma'(t)=\gamma_\infty(t)
+2\Rea\frac{\zeta'}{\zeta}(\sigma+\ii t),\quad
\gamma_\infty(t)=\Rea\psi(1/4+\ii t/2)-\log\pi,
\end{equation}
with removable values used on exceptional lines.
Write
\[
\kappa_F(x)=\int\overline{F(y)}F(y+x)\dd y,\qquad
\Gamma[F]=\int\gamma_\infty(t)|\widehat F(t)|^2\frac{\dd t}{2\pi},
\]
and, for any real \(\sigma\),
\begin{equation}\label{eq:W}
W_\sigma[F]=\sum_{p,m\ge1}(\log p)p^{-m\sigma}
\{\kappa_F(m\log p)+\kappa_F(-m\log p)\}.
\end{equation}
If the support of \(F\) has diameter at most \(L\), only delays
\(m\log p<L\) contribute. All active prime powers are retained.

\begin{theorem}[Boundary-resolvent representation]\label{thm:boundary}
For \(\sigma>1\) and \(F\in C_c^\infty(\mathbb R)\),
\begin{equation}\label{eq:boundary-formula}
B_\sigma[F]=\Gamma[F]-W_\sigma[F]+K_\sigma[F],
\end{equation}
where the correction is independently represented by
\begin{equation}\label{eq:short-boundary}
K_\sigma[F]=
2\Rea\Tr_{\chi\HH}
\left(C_\sigma(\Id_\chi-C_\sigma^2)^{-1}T_\sigma X_{F,e}\right).
\end{equation}
The kernel of \(X_{F,e}\) is
\(1_{x>0>y}\check a_e(x-y)\); it vanishes unless
\(-L<y<0<x<L\) and \(x-y<L\).
With \(c_\sigma=\norm{C_\sigma}<1\), the series
\begin{equation}\label{eq:return-series}
K_\sigma[F]=2\Rea\sum_{n\ge0}
\Tr_{\chi\HH}(C_\sigma^{2n+1}T_\sigma X_{F,e})
\end{equation}
has truncation error after \(n=N\) bounded by
\begin{equation}\label{eq:return-series-error}
\frac{2c_\sigma^{2N+3}}{\sqrt{1-c_\sigma^2}}\,
\norm{X_{F,e}}_{\Sone}.
\end{equation}
\end{theorem}
\begin{proof}
For fixed \(\sigma>1\), derivatives of the zeta ratio are bounded
at each order; the gamma ratio and its derivatives have at most
polynomial growth. Thus \eqref{eq:phase-comm} applies to every Schwartz
multiplier. The trace calculation of Lemma~\ref{lem:phase-trace}
extends from compact frequency tests by trace-norm approximation.
The absolutely convergent differentiated Euler series in
\eqref{eq:phase-deriv} gives the bulk
\(\Gamma[F]-W_\sigma[F]\).

For the remaining algebra suppress \(\sigma\), and set
\(s=(\Id_\chi-C^2)^{1/2}\).
Equation~\eqref{eq:projection} gives
\[
\Pi=L-H,\qquad H=T^*C^2s^{-2}T\ge0.
\]
Since \(0\le H\le L\), its source sandwich is trace class.
The first boundary formula is consequently
\begin{equation}\label{eq:long-boundary}
K_\sigma[F]=
\Tr(C_FC^2C_F^*)-\Tr(C_FHC_F^*)
+2\Rea\Tr_{\chi\HH}(CTPa(D)\chi).
\end{equation}
It follows directly from \eqref{eq:Qblocks}, without defining \(K\)
by arithmetic subtraction.

Real structure makes the frequency measure of \(\Pi\) even, and the
bulk derivative is even. We may therefore replace \(a\) by \(a_e\)
in \(K\). The multiplier \(a_e(D)\) commutes with \(\Fcal\).
The isometry \(J=T^*s^{-1}:\chi\HH\to(P-\Pi)\HH\)
gives on \(J\chi\HH\oplus\chi\HH\)
\[
\Fcal=\begin{pmatrix}-C&s\\s&C\end{pmatrix},
\qquad
\Pi+\Qcal=\begin{pmatrix}-C^2&sC\\Cs&C^2\end{pmatrix}.
\]
This subspace reduces \(\Fcal\): \(\Ran\Pi\) is itself invariant
under the involution. If the compression of \(a_e(D)\) is
\(\begin{pmatrix}A&B\\B^*&D_0\end{pmatrix}\), commutation gives
\[
sD_0-As=CB+BC,\qquad B=s^{-1}T X_{F,e}\in\Sone.
\]
The diagonal products \(AC^2,D_0C^2\) are trace class: this follows
by compressing \(a_e(D)\Qcal\), whose diagonal blocks are
\(-AC^2+BCs\) and \(B^*sC+D_0C^2\).
Multiplying the commutation relation on the right by \(Cs^{-1}C\),
and using bounded-factor cyclicity, gives
\[
\Tr(C^2D_0-C^2A)=2\Tr(C^3s^{-1}B).
\]
Adding the crossing traces reduces \eqref{eq:long-boundary}
to \(2\Rea\Tr(Cs^{-1}B)\), which is \eqref{eq:short-boundary}.

The Neumann series for \(s^{-2}\) proves
\eqref{eq:return-series}. The omitted bounded factor is
\(C^{2N+3}(\Id_\chi-C^2)^{-1}T\).
Multiplication by its adjoint and \eqref{eq:tt} give squared norm
at most \(c_\sigma^{4N+6}/(1-c_\sigma^2)\).
The trace-ideal inequality proves \eqref{eq:return-series-error}.
\end{proof}

\begin{remark}[Archimedean calibration]\label{rem:calibration}
As \(\sigma\to\infty\), \(D_\sigma,D_\sigma^{-1}\to\Id\) in norm.
The gap and the fixed trace-class crossing block show directly that
\eqref{eq:short-boundary} tends to its archimedean value.
If \(C_\infty\xi_n=\lambda_n\xi_n\) and
\(\zeta_n=T_\infty^*\xi_n/\sqrt{1-\lambda_n^2}\), that value is
\[
2\Rea\sum_n\frac{\lambda_n}{\sqrt{1-\lambda_n^2}}
\inner{\zeta_n}{a_e(D)\xi_n}.
\]
The smooth finite-interval cosine kernel makes \(C_\infty\) trace
class, so the sum is absolutely convergent.
This is the prolate correction of
\cite[Theorem 7]{CCWeil} in the present conventions.
The correction is generally nonzero.
\end{remark}

\section{Abel regularization without a cutoff gap}\label{sec:abel}

The norm gap in Proposition~\ref{prop:gap} is only established for
\(\sigma>1\). Bounded spectral calculus gives a replacement on every
line \(\sigma>1/2\).

\begin{proposition}[Positive Abel approximants]\label{prop:abel}
For \(0\le r<1\), define
\begin{align}
\Pi_{\sigma,r}
&=P-T_\sigma^*(\Id_\chi-rC_\sigma^2)^{-1}T_\sigma,\label{eq:abel}\\
H_{\sigma,r}
&=rT_\sigma^*C_\sigma^2(\Id_\chi-rC_\sigma^2)^{-1}T_\sigma.
\label{eq:abel-return}
\end{align}
Then
\[
0\le\Pi_\sigma\le\Pi_{\sigma,r}\le L_\sigma,\qquad
\Pi_{\sigma,r}\downarrow\Pi_\sigma,\qquad
H_{\sigma,r}\uparrow H_\sigma
\]
in the strong operator topology as \(r\uparrow1\), and
\(\Pi_{\sigma,r}=L_\sigma-H_{\sigma,r}\).
For every compact smooth frequency multiplier \(b\), these monotone
limits pass through the positive smoothed traces.
\end{proposition}
\begin{proof}
Put \(A=T_\sigma^*T_\sigma\) on \(P\HH\).
Using \(T_\sigma T_\sigma^*=\Id_\chi-C_\sigma^2\),
functional calculus gives
\[
\Pi_{\sigma,r}
=\frac{(1-r)(\Id_{P\HH}-A)}
{(1-r)\Id_{P\HH}+rA}.
\]
At a spectral value \(\alpha\in[0,1]\) the multiplier is
\((1-r)(1-\alpha)/(1-r+r\alpha)\). It equals one at
\(\alpha=0\), decreases to zero at every \(\alpha>0\), and lies
between zero and \(1-\alpha\). This proves all operator assertions.
Lemma~\ref{lem:phase-trace} supplies the finite dominating trace
\(\Tr(b(D)L_\sigma b(D)^*)\), so monotone convergence applies.
\end{proof}

In particular \(\nu_\sigma\) is locally finite for every
\(\sigma>1/2\). A gap is unnecessary for this assertion.
The approximants in \eqref{eq:abel} are generally not projections.
Their positive sandwich trace is
\(\Tr(b(D)\Pi_{\sigma,r}b(D)^*)\); it must not be replaced by
\(\norm{b(D)\Pi_{\sigma,r}}_{\Stwo}^2\).

\section{Uniform frequency bounds}\label{sec:tails}

We now prove \eqref{eq:main-window} and extend local trace formulas to
all Schwartz source multipliers. The estimates are uniform when a zero
or the zeta pole approaches the chosen vertical line.

\subsection{Local factorization and a uniform energy bound}

The classical inputs are the local zero count and local partial-fraction
formula for \(\zeta'/\zeta\); see
\cite[Propositions 16 and 19]{Tao} and \cite{Elkies}.
Counting multiplicities,
\begin{equation}\label{eq:zero-count}
\#\{\rho=\beta+\ii\gamma:|\gamma-j|\le5\}
=O(\log(2+|j|)).
\end{equation}
Reflection in the critical line extends a count in a strip bounded
away from zero to the full critical strip.
For fixed \(S>1/2\), \(1/2\le\sigma\le S\), and
\(t\in J_j=(j-3,j+3)\), the local logarithmic-derivative estimate gives
\begin{equation}\label{eq:local-log}
\frac{\zeta'}{\zeta}(\sigma+\ii t)
-\sum_{|\gamma-j|\le5}\frac{m_\rho}{\sigma+\ii t-\rho}
+\frac1{\sigma+\ii t-1}
=O_S(\log(2+|j|)),
\end{equation}
after removal of the displayed singularities.
To obtain this version, first extract zeros in a fixed small disk
about \(\sigma+\ii t\). Enlarging the extracted set to the
displayed list adds only \(O(\log(2+|j|))\) terms whose denominators
are bounded away from zero. The estimate for bounded \(j\) follows
by factoring finitely many singularities on a compact set.

For \(d\ne0\), define
\[
B_{d,\gamma}(t)=\frac{t-\gamma-\ii d}{t-\gamma+\ii d},
\qquad B_{0,\gamma}=1.
\]
The zero-width convention concerns the phase on the line itself.
The phase derivative of a nonconstant factor is
\begin{equation}\label{eq:Bderiv}
\frac{2d}{d^2+(t-\gamma)^2}.
\end{equation}
Include the inverse pole factor
\[
R_\sigma(t)=\frac{t+\ii(\sigma-1)}{t-\ii(\sigma-1)}
\quad(\sigma\ne1),\qquad R_1=1,
\]
and, on \(J_j\), set
\begin{equation}\label{eq:local-product}
\mathcal B_{j,\sigma}
=R_\sigma\prod_{|\gamma-j|\le5}
B_{\sigma-\beta,\gamma}^{\,m_\rho},
\qquad w_{j,\sigma}=v_\sigma/\mathcal B_{j,\sigma}.
\end{equation}
All factors have modulus one. Equation~\eqref{eq:local-log}
and the \(O(\log(2+|t|))\) bound for the archimedean phase derivative
show that the removable extension of \(w_{j,\sigma}\) obeys
\begin{equation}\label{eq:background}
\sup_{J_j}|w_{j,\sigma}'|
\le C_S\log(2+|j|).
\end{equation}
The estimate includes \(\sigma=1\) and lines containing zeros.
The real part of \(1/(\ii(t-\gamma))\) is zero away from
\(\gamma\); its phase contribution on the exceptional line is
removed, not interpreted as a delta mass.

For a scalar function \(q\), write
\[
\mathcal E(q)=\iint
\frac{|q(t)-q(s)|^2}{(t-s)^2}\dd t\dd s.
\]

\begin{lemma}[Width-independent factor energy]\label{lem:energy}
For \(d\ne0\), \(\mathcal E(B_{d,\gamma})=4\pi^2\).
There is a fixed real cutoff \(b_0\in C_c^\infty((-2,2))\),
equal to one on \([0,1]\), such that, with \(b_j(t)=b_0(t-j)\),
\begin{equation}\label{eq:localized-energy}
\mathcal E(b_jv_\sigma)\le C_S\log^2(2+|j|)
\end{equation}
uniformly for \(1/2<\sigma\le S\).
\end{lemma}
\begin{proof}
Direct subtraction gives
\[
\frac{|B_{d,\gamma}(t)-B_{d,\gamma}(s)|^2}{(t-s)^2}
=\frac{4d^2}
{((t-\gamma)^2+d^2)((s-\gamma)^2+d^2)}.
\]
Integration proves the first assertion, independently of the sign
or magnitude of \(d\). The inverse pole factor has the same energy.

Extend \(w_{j,\sigma}\) constantly outside \(J_j\).
This extension is continuous and Lipschitz. Let
\(f_{j,\sigma}=b_jw_{j,\sigma}\). The translation estimates
\[
\norm{f(\cdot+h)-f}_2\le |h|\norm{f'}_2\quad (|h|\le1),
\qquad
\norm{f(\cdot+h)-f}_2\le2\norm f_2\quad (|h|>1)
\]
imply
\[
\mathcal E(f_{j,\sigma})
\le C\bigl(\norm{f_{j,\sigma}}_2^2+
\norm{f_{j,\sigma}'}_2^2\bigr)
\le C_S\log^2(2+|j|).
\]
For a product \(B=\prod_{k=1}^{M}B_k\) of unit-modulus factors,
telescoping and Cauchy--Schwarz give
\[
|B(t)-B(s)|^2\le M\sum_{k=1}^{M}|B_k(t)-B_k(s)|^2.
\]
Here \(M=O(\log(2+|j|))\), including the pole factor.
Using
\[
|f(t)B(t)-f(s)B(s)|^2
\le2|f(t)-f(s)|^2+
2\norm f_\infty^2|B(t)-B(s)|^2
\]
proves \eqref{eq:localized-energy}. The identity
\(b_jv_\sigma=f_{j,\sigma}\mathcal B_{j,\sigma}\) holds globally
because \(b_j\) is supported strictly inside \(J_j\).
No global Blaschke factorization has been assumed.
\end{proof}

\subsection{Positive traces and source tails}

\begin{proposition}[Uniform trace control]\label{prop:tails}
The estimate \eqref{eq:main-window} holds. For every
\(f\in\Sch(\mathbb R)\) and any of the three positive measures there,
\begin{equation}\label{eq:tail-sum}
\int |f(t)|^2\dd\nu_G(t)
\le M_S\sum_{j\in\mathbb Z}
\log^2(2+|j|)\sup_{t\in[j,j+1]}|f(t)|^2.
\end{equation}
For each integer \(N>1\), its tail on \(|t|>A\) is
\(O_{S,f,N}(A^{1-2N}\log^2(2+A))\), uniformly in \(\sigma\).
\end{proposition}
\begin{proof}
Factor \(L_\sigma=A_\sigma^*A_\sigma\), with
\(A_\sigma=\chi V_\sigma P\), where \(V_\sigma=v_\sigma(D)\).
Then
\[
A_\sigma b_j(D)^*
=\chi(v_\sigma\overline b_j)(D)P+
\chi V_\sigma[P,b_j(D)^*].
\]
Lemma~\ref{lem:hankel} and \eqref{eq:localized-energy} imply
\[
\Tr(b_j(D)L_\sigma b_j(D)^*)
=\norm{A_\sigma b_j(D)^*}_{\Stwo}^2
\le C_S\log^2(2+|j|).
\]
Since \(b_j=1\) on \([j,j+1]\), this proves the estimate for
\(\nu_{L_\sigma}\). Positive domination gives it for \(\nu_\sigma\).
Using \eqref{eq:Cfactor}, factor the small-cutoff sandwich through
\(PV_\sigma\chi b_j(D)^*\), and commute \(b_j(D)^*\) past
\(\chi\). Exactly the same estimates prove the bound for
\(\nu_{C_\sigma^2}\). This includes any subspace at spectral
value one of \(C_\sigma^2\).

Summing the unit-window estimates gives \eqref{eq:tail-sum}.
For the tail, use the order-\(N\) Schwartz bound and sum over
\(|j|>A-2\). The resulting scalar series has the asserted bound.
\end{proof}

The same bounds hold for \(H_\sigma\) and all \(\Pi_{\sigma,r}\),
by positive domination. Thus full compact smooth position sources
require no auxiliary finite-rank cutoff.

\begin{proposition}[Full-source boundary identity]\label{prop:full-boundary}
For every \(f\in\Sch(\mathbb R)\) and \(\sigma>1/2\),
\begin{align}
\Tr(f(D)\Pi_\sigma f(D)^*)
&=\int |f(t)|^2\phi_\sigma'(t)\frac{\dd t}{2\pi}+K_\sigma[f],
\label{eq:full-boundary}\\
K_\sigma[f]
&=\Tr(f(D)C_\sigma^2 f(D)^*)
-\Tr(f(D)H_\sigma f(D)^*)\notag\\
&\hspace{1em}+2\Rea\Tr_{\chi\HH}
\bigl(C_\sigma T_\sigma P|f|^2(D)\chi\bigr).
\label{eq:full-K}
\end{align}
All displayed quantities are finite. The two positive traces have
uniform frequency tails on bounded \(\sigma\)-strips.
\end{proposition}
\begin{proof}
Local factorization also gives
\begin{equation}\label{eq:phase-L1}
\int_j^{j+1}|\phi_\sigma'(t)|\dd t
\le C_S\log(2+|j|).
\end{equation}
Indeed each rational factor has total absolute phase variation at
most \(2\pi\), and the remainder satisfies \eqref{eq:background}.
This proves absolute convergence and uniform Schwartz tails for
the bulk integral.

Apply Lemma~\ref{lem:phase-trace} and \(\Pi_\sigma=L_\sigma-H_\sigma\)
to \(f_N=f b_0(t/N)\), replacing \(b_0\) if necessary by a
cutoff equal to one near zero. Then \(f_N\to f\) in the Schwartz
topology. Proposition~\ref{prop:tails} treats the positive trace
terms. Lemma~\ref{lem:hankel} gives
\[
P|f_N|^2(D)\chi\longrightarrow P|f|^2(D)\chi
\quad\hbox{in trace norm}.
\]
The remaining factor \(C_\sigma T_\sigma\) is bounded in norm by one,
so the crossing traces converge. Equation~\eqref{eq:phase-L1}
treats the bulk integral and proves the result.
This proof assigns the crossing scalar directly to a trace-class
cutoff product; it does not require a global trace-class assertion
for an unsandwiched phase operator.
\end{proof}

\section{Localization at the critical boundary}\label{sec:critical}

\subsection{Strong collapse and local phase factors}

The functional equation, in the usual Riemann--Siegel convention
\cite{DLMF}, gives \(v_{1/2}(t)=1\) away from zero ordinates.
Bounded pointwise convergence therefore implies
\(V_\sigma\to\Id\) strongly as \(\sigma\downarrow1/2\).
Equations~\eqref{eq:Lfactor}--\eqref{eq:Cfactor} give
\begin{equation}\label{eq:strong}
\Fcal_\sigma\longrightarrow\J,\qquad
L_\sigma\longrightarrow0,\qquad
C_\sigma\longrightarrow0,\qquad
\Pi_\sigma\longrightarrow0
\quad\hbox{strongly}.
\end{equation}
For the last assertion,
\(\norm{\Pi_\sigma h}^2\le\inner h{L_\sigma h}\).
In particular, for fixed finite-rank \(E\) and bounded \(f(D)\),
\(\norm{f(D)\Pi_\sigma E}_{\Stwo}\to0\).
A fixed compact frequency multiplier is generally infinite rank,
so this observation does not determine its trace limit.

Let \(\epsilon=\sigma-1/2>0\). On a fixed compact frequency interval,
choose a slightly larger rectangle about the critical line with
boundary free of zeros and no off-line zero in its interior.
There are only finitely many zeros in a bounded rectangle, so its
horizontal width can be chosen this way. Local analytic factorization
gives, on a neighborhood of the interval,
\begin{equation}\label{eq:critical-factor}
v_{1/2+\epsilon}(t)=B_\epsilon(t)w_\epsilon(t),
\qquad
B_\epsilon(t)=\prod_\gamma
\left(\frac{t-\gamma-\ii\epsilon}
{t-\gamma+\ii\epsilon}\right)^{m_\gamma},
\qquad
w_\epsilon\longrightarrow1\quad\hbox{in }C^\infty.
\end{equation}
The product includes every critical zero in the rectangle. The factored
remainder is smooth jointly in \(\epsilon,t\); the functional equation
identifies its endpoint value. Zeros outside the rectangle do not
require any location hypothesis.

Every pole of \(B_\epsilon\) lies in the lower frequency half-plane.
Its inverse Fourier kernel is supported on \((-\infty,0]\).
For example
\begin{equation}\label{eq:onefactor-kernel}
\mathcal F^{-1}\left(\frac{t-\ii h}{t+\ii h}\right)(x)
=\delta_0(x)-2h e^{hx}1_{x<0},\qquad h>0.
\end{equation}
Partial fractions or convolution proves the corresponding support
statement for the finite product, and hence
\begin{equation}\label{eq:triangular-B}
P B_\epsilon(D)\chi=0.
\end{equation}

\begin{lemma}[Local phase concentration]\label{lem:phase-limit}
As \(\sigma\downarrow1/2\),
\[
\phi_\sigma'(t)\frac{\dd t}{2\pi}\longrightarrow\mucrit
\]
on compact smooth tests and on all Schwartz tests.
\end{lemma}
\begin{proof}
Differentiating the local factorization \eqref{eq:critical-factor}
gives a smooth term tending uniformly to zero and, for each critical
zero of multiplicity \(m_\gamma\), the term
\[
\frac{2m_\gamma\epsilon}{\epsilon^2+(t-\gamma)^2}.
\]
After division by \(2\pi\), this tends to
\(m_\gamma\delta_\gamma\). The local conclusion follows from the
Poisson-kernel approximate identity. The unit-window variation bound
\eqref{eq:phase-L1} and the local zero count give uniform Schwartz
tails for the prelimit and limit measures, respectively.
This proves the global assertion.
\end{proof}

\subsection{Small-cutoff and crossing terms}

\begin{lemma}[Vanishing localized cutoff]\label{lem:Csmall}
For every \(b\in C_c^\infty(\mathbb R)\),
\begin{equation}\label{eq:Csmall}
\norm{C_\sigma b(D)}_{\Stwo}\longrightarrow0.
\end{equation}
Furthermore,
\begin{equation}\label{eq:crosssmall}
\Tr_{\chi\HH}\bigl(C_\sigma T_\sigma P|f|^2(D)\chi\bigr)
\longrightarrow0
\quad\hbox{for every } f\in\Sch(\mathbb R).
\end{equation}
\end{lemma}
\begin{proof}
Choose the factorization \eqref{eq:critical-factor} on a
neighborhood of \(\operatorname{supp}b\), and set
\(q_\epsilon=b(w_\epsilon-1)\), extended by zero.
Then \(q_\epsilon\to0\) in the Schwartz topology.
Using \(C_\sigma=\J P V_\sigma\chi\),
\(\chi b(D)=b(D)\chi-[b(D),\chi]\), and
\eqref{eq:triangular-B}, one obtains the exact identity
\begin{align}\label{eq:Csmall-id}
C_\sigma b(D)=\J\{&
Pq_\epsilon(D)\chi B_\epsilon(D)\chi\\
&+P(B_\epsilon(D)-V_\sigma)[b(D),\chi]\}.\notag
\end{align}
The first term tends to zero in Hilbert--Schmidt norm by
Lemma~\ref{lem:hankel}. Both \(B_\epsilon(D)\) and \(V_\sigma\)
converge strongly to the identity. Thus the second term is a
uniformly bounded strongly null family times a fixed
Hilbert--Schmidt operator, and tends to zero in that ideal.
This last implication follows by finite-rank approximation.

For \eqref{eq:crosssmall}, \(P|f|^2(D)\chi\) is trace class,
whereas \(C_\sigma T_\sigma\to0\) strongly by
\eqref{eq:strong}. Finite-rank approximation in trace norm
proves the scalar limit.
\end{proof}

Combining these lemmas with \eqref{eq:L-trace} gives
\begin{equation}\label{eq:local-Llimit}
\Tr(b(D)L_\sigma b(D)^*)\longrightarrow
\int|b|^2\dd\mucrit
\quad (b\in C_c^\infty).
\end{equation}
The uniform tails of Proposition~\ref{prop:tails} extend this
limit to every Schwartz \(f\). They likewise extend
\(\Tr(b(D)C_\sigma^2b(D)^*)\to0\) to all such \(f\).
Thus \eqref{eq:full-K} implies
\begin{align}
\Tr(f(D)\Pi_\sigma f(D)^*)
&=\int|f|^2\dd\mucrit-\Tr(f(D)H_\sigma f(D)^*)+o_f(1),
\label{eq:deficit}\\
K_\sigma[f]&=-\Tr(f(D)H_\sigma f(D)^*)+o_f(1).
\label{eq:Kdeficit}
\end{align}

\begin{corollary}[Possible exact-projection limits]\label{cor:subseq}
Every sequence \(\sigma_n\downarrow1/2\) has a vaguely convergent
subsequence of \(\nu_{\sigma_n}\). Any limit is bounded above by
\(\mucrit\) and therefore has precisely the form stated in
Theorem~\ref{thm:main}.
\end{corollary}
\begin{proof}
The unit-window bound gives compactness of the positive measures
on each bounded interval. A diagonal subsequence gives vague
convergence on \(\mathbb R\). Since
\(\nu_{\sigma}\le\nu_{L_\sigma}\), \eqref{eq:local-Llimit}
implies \(\nu_*\le\mucrit\).
One may pass from squared smooth tests to all compact continuous
tests by choosing a smooth function equal to one on their support
and uniformly approximating there; the local mass bounds justify
this step. A positive measure dominated by the discrete
\(\mucrit\) has no other support and has atom weights bounded
by its multiplicities.
\end{proof}

\section{The return defect at cutoff spectral value one}\label{sec:return}

Let
\begin{equation}\label{eq:polar}
T_\sigma=(\Id_\chi-C_\sigma^2)^{1/2}W_\sigma^{\mathrm{pol}}
\end{equation}
be the polar decomposition in this orientation.
The superscript distinguishes the partial isometry from the prime sum
\(W_\sigma[F]\). Its initial space is
\((\ker T_\sigma)^\perp\subset P\HH\), and its final space is
\(\ker(\Id_\chi-C_\sigma^2)^\perp\).
Consequently
\begin{equation}\label{eq:Hpolar}
H_\sigma=(W_\sigma^{\mathrm{pol}})^*C_\sigma^2
W_\sigma^{\mathrm{pol}}.
\end{equation}
For \(b\in C_c^\infty(\mathbb R)\), define the positive spectral measure
\begin{equation}\label{eq:eta}
\eta_{\sigma,b}(E)=
\Tr\!\left(b(D)(W_\sigma^{\mathrm{pol}})^*C_\sigma^2
1_E(C_\sigma^2)W_\sigma^{\mathrm{pol}}b(D)^*\right),
\quad E\subset[0,1].
\end{equation}
It is finite, its total mass is \(R_\sigma[b]\), and
\(\eta_{\sigma,b}(\{1\})=0\) at each fixed \(\sigma\).
The last fact follows from the final space of the polar partial
isometry. It gives no uniform exclusion of mass close to one.

\begin{theorem}[Return localization and Abel error]\label{thm:return}
For \(b\in C_c^\infty(\mathbb R)\) and \(0<\delta<1\),
\begin{equation}\label{eq:return-localize}
\eta_{\sigma,b}([0,1-\delta])
\le \delta^{-1}\norm{C_\sigma T_\sigma b(D)^*}_{\Stwo}^2
\longrightarrow0
\quad(\sigma\downarrow1/2).
\end{equation}
For \(0\le r<1\),
\begin{align}
\Tr\!\left(b(D)(\Pi_{\sigma,r}-\Pi_\sigma)b(D)^*\right)
&=\int_{[0,1]}\frac{1-r}{1-r\lambda}\dd\eta_{\sigma,b}(\lambda)
\label{eq:Abel-error}\\
&\le
\frac{1-r}{1-r+r\delta}\eta_{\sigma,b}([0,1])
+\eta_{\sigma,b}((1-\delta,1)).
\notag
\end{align}
At each fixed \(\sigma\), this error tends to zero as \(r\uparrow1\).
For each fixed \(r<1\),
\begin{equation}\label{eq:fixed-Abel}
\Tr(f(D)\Pi_{\sigma,r}f(D)^*)\longrightarrow
\int |f|^2\dd\mucrit
\quad\hbox{for every }f\in\Sch(\mathbb R).
\end{equation}
\end{theorem}
\begin{proof}
Set \(b^\sharp(t)=\overline{b(-t)}\).
Since \(\Fcal_\sigma b(D)^*=b^\sharp(D)\Fcal_\sigma\),
\begin{equation}\label{eq:CTsmall}
C_\sigma T_\sigma b(D)^*
=C_\sigma b^\sharp(D)\Fcal_\sigma P
+C_\sigma\Fcal_\sigma[P,b(D)^*].
\end{equation}
The first term tends to zero in Hilbert--Schmidt norm by
Lemma~\ref{lem:Csmall}. The second does also: its commutator is
fixed Hilbert--Schmidt, and
\(C_\sigma\Fcal_\sigma\to0\) strongly.
This proves the limit on the right of \eqref{eq:return-localize}.

Under \eqref{eq:polar}, the squared norm in that equation contains
the spectral weight \(\lambda(1-\lambda)\), while \(\eta\)
contains \(\lambda\). On \([0,1-\delta]\), the latter is at most
\(\delta^{-1}\lambda(1-\lambda)\). This proves the inequality.

Subtracting the exact kernel projection from \eqref{eq:abel}
and using the polar decomposition gives
\[
\Pi_{\sigma,r}-\Pi_\sigma
=(W_\sigma^{\mathrm{pol}})^*
\frac{(1-r)C_\sigma^2}{\Id_\chi-rC_\sigma^2}
W_\sigma^{\mathrm{pol}}.
\]
This proves \eqref{eq:Abel-error}. Splitting its spectral integral
at \(1-\delta\) proves the bound. Dominated convergence applies
at fixed \(\sigma\), since there is no atom at one.

For fixed \(r<1\), \eqref{eq:abel-return} gives
\[
0\le\Tr(b(D)H_{\sigma,r}b(D)^*)
\le\frac{r}{1-r}
\norm{C_\sigma T_\sigma b(D)^*}_{\Stwo}^2\longrightarrow0.
\]
Combine this with \(\Pi_{\sigma,r}=L_\sigma-H_{\sigma,r}\)
and \eqref{eq:local-Llimit}. This proves \eqref{eq:fixed-Abel}
for compact smooth \(b\). The domination
\(0\le\Pi_{\sigma,r}\le L_\sigma\) and the uniform source tails
extend it to every Schwartz \(f\).
\end{proof}

Theorems~\ref{thm:return} and~\ref{thm:main} leave open precisely
the total return mass. For a fixed compact frequency test, the
sufficient condition
\begin{equation}\label{eq:uniform-integrability}
\lim_{\delta\downarrow0}\,
\limsup_{\sigma\downarrow1/2}
\eta_{\sigma,b}((1-\delta,1))=0
\end{equation}
would imply \(R_\sigma[b]\to0\), using
\eqref{eq:return-localize}. Conversely, vanishing total mass implies
\eqref{eq:uniform-integrability}. A mere positive gap for each
prelimit parameter does not supply this condition.

For a jointly chosen \(r=r(\sigma)\uparrow1\), the sufficient
estimate
\[
\frac{\norm{C_\sigma T_\sigma b(D)^*}_{\Stwo}^2}
{1-r(\sigma)}\longrightarrow0
\]
gives the critical-zero limit for the corresponding Abel trace.
This is a source-dependent sufficient criterion, not a proved
uniform rate for the zeta family.

\section{A rational family with complete projection loss}\label{sec:model}

The following example concerns a different global phase. Its role
is to test what can be inferred from local critical factorization,
strong convergence, and a positive prelimit gap.

\begin{theorem}[Exact rational model]\label{thm:model}
For \(0<\epsilon<1\), set \(R=\epsilon^{-1}\) and
\begin{equation}\label{eq:model-phase}
B_h(t)=\frac{t-\ii h}{t+\ii h},\qquad
v_\epsilon^{\mathrm{mod}}(t)=-\frac{B_\epsilon(t)}{B_R(t)},
\qquad
q=\frac{R-\epsilon}{R+\epsilon}
=\frac{1-\epsilon^2}{1+\epsilon^2}.
\end{equation}
Define the cutoff operators as before from
\(\Fcal_\epsilon^{\mathrm{mod}}=\J v_\epsilon^{\mathrm{mod}}(D)\).
Then the following statements hold:
\begin{enumerate}
\item \(v_\epsilon^{\mathrm{mod}}=B_\epsilon w_\epsilon\)
locally, with \(w_\epsilon=-B_R^{-1}\to1\) in \(C^\infty\);
the phase measures tend locally to \(\delta_0\).
\item The cutoff gap is strictly positive for every \(\epsilon\),
\[
1-\norm{C_\epsilon^{\mathrm{mod}}}^2
=1-q^2=\frac{4\epsilon^2}{(1+\epsilon^2)^2}.
\]
Nevertheless \(\Pi_\epsilon^{\mathrm{mod}}=0\) exactly.
\item For every Schwartz \(b\), the majorant trace tends to
\(|b(0)|^2\), and the entire limiting mass lies in the return defect.
\end{enumerate}
\end{theorem}
\begin{proof}
The phase derivative is
\begin{equation}\label{eq:model-derivative}
(\phi_\epsilon^{\mathrm{mod}})'(t)
=\frac{2\epsilon}{t^2+\epsilon^2}
-\frac{2R}{t^2+R^2}.
\end{equation}
On compact tests, the first term divided by \(2\pi\) tends
to \(\delta_0\), and the second tends to zero. The asserted
factorization and smooth background limit follow directly.

Partial fractions, with the Fourier convention \eqref{eq:fourier},
give
\begin{equation}\label{eq:model-kernel}
\check v_\epsilon^{\mathrm{mod}}(x)
=-\delta_0(x)-2\epsilon q e^{\epsilon x}1_{x<0}
+2Rq e^{-Rx}1_{x>0}.
\end{equation}
Let
\[
u_\epsilon(x)=\sqrt{2\epsilon}\,e^{-\epsilon x}1_{x>0},
\qquad
\xi_R(x)=\sqrt{2R}\,e^{Rx}1_{x<0}.
\]
These are unit vectors. The kernel of \(\Fcal_\epsilon^{\mathrm{mod}}\)
is \(\check v_\epsilon^{\mathrm{mod}}(-x-y)\), so
\begin{equation}\label{eq:model-blocks}
P\Fcal_\epsilon^{\mathrm{mod}}P
=-q\,|u_\epsilon\rangle\langle u_\epsilon|,
\qquad
C_\epsilon^{\mathrm{mod}}
=q\,|\xi_R\rangle\langle\xi_R|.
\end{equation}
Therefore
\[
L_\epsilon^{\mathrm{mod}}
=q^2|u_\epsilon\rangle\langle u_\epsilon|,\qquad
\norm{L_\epsilon^{\mathrm{mod}}}=q^2<1.
\]
An orthogonal projection dominated by a positive operator of
norm less than one is zero. Since \(\Pi\le L\), this proves
\(\Pi_\epsilon^{\mathrm{mod}}=0\) and \(H_\epsilon^{\mathrm{mod}}=L_\epsilon^{\mathrm{mod}}\).

Finally,
\begin{equation}\label{eq:model-source}
\norm{b(D)u_\epsilon}_2^2
=\int |b(t)|^2
\frac{2\epsilon}{\epsilon^2+t^2}\frac{\dd t}{2\pi}
\longrightarrow |b(0)|^2.
\end{equation}
This proves the trace assertion for Schwartz \(b\).
\end{proof}

The spectral defect and both orders of limits are explicit.
Indeed \(T_\epsilon^{\mathrm{mod}}u_\epsilon
=-\sqrt{1-q^2}\,\xi_R\), so
\begin{equation}\label{eq:model-eta}
\eta_{\epsilon,b}^{\mathrm{mod}}
=q^2\norm{b(D)u_\epsilon}_2^2\,\delta_{q^2}.
\end{equation}
It has no atom at one for any \(\epsilon>0\), but asymptotically
all its mass approaches one.

\begin{corollary}[Abel crossover]\label{cor:crossover}
In the rational model,
\begin{equation}\label{eq:model-Abel}
\Pi_{\epsilon,r}^{\mathrm{mod}}
=\frac{q^2(1-r)}{1-rq^2}
|u_\epsilon\rangle\langle u_\epsilon|.
\end{equation}
Writing \(d_\epsilon=1-r(\epsilon)\), its smoothed trace tends to
\[
\begin{cases}
|b(0)|^2,& d_\epsilon/\epsilon^2\to\infty,\\
0,& d_\epsilon/\epsilon^2\to0,\\
\dfrac{c}{c+4}|b(0)|^2,& d_\epsilon/\epsilon^2\to c\in(0,\infty).
\end{cases}
\]
\end{corollary}
\begin{proof}
On the span of \(u_\epsilon\), \(T^*T\) has eigenvalue
\(1-q^2\); on its orthogonal complement it has eigenvalue one.
The functional-calculus formula in Proposition~\ref{prop:abel}
gives \eqref{eq:model-Abel}.
Use \(1-q^2\sim4\epsilon^2\) and \eqref{eq:model-source}.
\end{proof}

Thus local phase concentration can be canceled completely by an
exact-projection return term. The example shares the local structure
used in Section~\ref{sec:critical}, but it is not a counterexample
about the actual zeta phase. A theorem determining the weights
\(c_\gamma\) for zeta must use additional information about that
global phase.

\section{Arithmetic comparison and crossing residues}\label{sec:arithmetic}

The operator results above concern the frequency measure on critical
zeros. We now retain the arithmetic terms needed to compare it with
the full explicit formula. The residue calculation is classical;
its purpose here is to fix the normalization and identify which
terms remain after the operator limit.

\subsection{Entire source weights}

Let \(F\in C_c^\infty(\mathbb R)\), possibly complex, and define
\begin{equation}\label{eq:entire-a}
a(z)=\widehat F(z)\overline{\widehat F(\overline z)}.
\end{equation}
It is entire, satisfies \(a(\overline z)=\overline{a(z)}\),
and equals \(|\widehat F(t)|^2\) on the real line.
It need not be even. For every \(D<\infty\) and every integer \(N\),
compact support and integration by parts give
\begin{equation}\label{eq:strip-decay}
\sup_{|\Ima z|\le D}(1+|\Rea z|)^N|a(z)|<\infty.
\end{equation}
The local zero count implies \(N(T)=O(T\log(T+2))\).
Consequently all zero sums below converge absolutely, uniformly
for bounded imaginary shifts in the source weight.

Define the phase bulk
\[
A_\sigma[F]=\int a(t)\phi_\sigma'(t)\frac{\dd t}{2\pi},
\]
and the crossing terms
\begin{align}
\mathcal P_\sigma[F]
&=
\begin{cases}
2\Rea a(\ii(1-\sigma)),&1/2<\sigma<1,\\
a(0),&\sigma=1,\\
0,&\sigma>1,
\end{cases}
\label{eq:Pcross}\\
\mathcal Z_\sigma[F]
&=2\Rea\sum_{\beta>\sigma}m_{\beta+\ii\gamma}
a\bigl(\gamma+\ii(\beta-\sigma)\bigr)
+\sum_{\beta=\sigma}m_{\beta+\ii\gamma}a(\gamma).
\label{eq:Zcross}
\end{align}
Both signs of \(\gamma\) occur. The on-line sum in
\eqref{eq:Zcross} is a half-residue contribution relative to
the strictly right-hand sum.

\begin{theorem}[Complete phase comparison]\label{thm:crossings}
For every \(\sigma>1/2\),
\begin{equation}\label{eq:crossing-formula}
A_\sigma[F]=\Gamma[F]-W_\sigma[F]
+\mathcal P_\sigma[F]-\mathcal Z_\sigma[F].
\end{equation}
The value on an exceptional line uses the removable phase itself.
Moreover,
\begin{equation}\label{eq:critical-explicit}
\Gamma[F]-W_{1/2}[F]
=\Zcrit[F]+\Zoff[F]-2\Rea a(\ii/2),
\end{equation}
where \(\Zoff\) is as in \eqref{eq:intro-arithmetic}.
\end{theorem}
\begin{proof}
The differentiated Euler series proves
\eqref{eq:crossing-formula} when \(\sigma>1\).
For \(1/2<\sigma<2\), initially excluding the pole line and real
parts of zeros, move the logarithmic-derivative contour to
\(\Rea s=2\). Consider
\begin{equation}\label{eq:contour-integrand}
a(-\ii(s-\sigma))\frac{\zeta'}{\zeta}(s)\frac{\dd s}{2\pi\ii}.
\end{equation}
Choose heights \(T_n\in[n,n+1]\) at distance at least
\(c/\log(n+2)\) from every nearby zero ordinate. The local zero
count constructs such heights by excluding intervals of sufficiently
small total length. Conjugation gives the same property at \(-T_n\).
The local logarithmic-derivative estimate then bounds
\(\zeta'/\zeta\) by \(O(\log^2(T_n+2))\) along the horizontal
sides, uniformly for real parts in \([\sigma,2]\).
Together with \eqref{eq:strip-decay}, this makes the horizontal
integrals vanish.

Orient the left side upward and the right side downward. The
orientation is clockwise, so the left integral equals the right
integral minus the residues. The logarithmic derivative has residue
\(-1\) at the pole and \(+m_\rho\) at a zero. The crossing terms
from \eqref{eq:contour-integrand} are therefore
\[
1_{\sigma<1}a(-\ii(1-\sigma))
-\sum_{\beta>\sigma}m_{\beta+\ii\gamma}
a\bigl(\gamma-\ii(\beta-\sigma)\bigr).
\]
On the right side the Euler series is absolutely convergent.
Shifting the entire source weight gives
\[
\int a(t-\ii(2-\sigma))e^{-\ii t\log n}\frac{\dd t}{2\pi}
=n^{2-\sigma}\kappa_F(-\log n).
\]
Thus its integral is
\(-\sum_{n\ge2}\Lambda(n)n^{-\sigma}\kappa_F(-\log n)\).
Interchange the integral and series on the line two first;
compact correlation support then makes the final sum finite.

Repeat the calculation with \(a(-z)\), and add the finite-height
identities before taking the limit. The left side becomes
\[
\int_{-T_n}^{T_n}a(t)
\left(\frac{\zeta'}{\zeta}(\sigma+\ii t)+
\frac{\zeta'}{\zeta}(\sigma-\ii t)\right)\frac{\dd t}{2\pi}.
\]
The paired integral is absolutely convergent by
\eqref{eq:phase-L1} and the archimedean bound.
No separate absolute convergence of the imaginary logarithmic
derivative is needed. The prime terms add to \(-W_\sigma[F]\).
Conjugation of the zero set and Schwarz reflection of \(a\)
put the residue terms in the form
\(\mathcal P_\sigma-\mathcal Z_\sigma\).
This proves \eqref{eq:crossing-formula} on nonexceptional lines.

At an exceptional line, take the average of the two neighboring
limits through nonexceptional lines. On a fixed bounded frequency
window, local factorization shows that a zero's two Cauchy-kernel
limits give opposite atoms; their average is its removable phase
value. The pole behaves in the same way with the opposite sign.
The uniform tails from \eqref{eq:phase-L1} extend this averaging
to the full source. The absolute and uniform convergence of the
shifted zero sums justifies averaging their residues. This produces
exactly the on-line term in \eqref{eq:Zcross} and
\(\mathcal P_1=a(0)\).

Finally Lemma~\ref{lem:phase-limit} gives
\(A_\sigma[F]\to\Zcrit[F]\). The prime sum is finite and
continuous in \(\sigma\). Dominated convergence in the crossing
sums gives
\(\mathcal Z_\sigma[F]\to\Zoff[F]\) and
\(\mathcal P_\sigma[F]\to2\Rea a(\ii/2)\).
At exceptional parameters approaching \(1/2\), the on-line
terms are controlled by the same uniform summable bound.
Taking the limit proves \eqref{eq:critical-explicit}.
\end{proof}

As checks on the signs, crossing the pole line from right to left
changes the bulk by \(+2a(0)\), whereas crossing a zero line changes
its local bulk contribution by \(-2m_\rho a(\gamma)\).
Preparation at the fixed pole moments
\(\widehat F(\pm\ii/2)=0\) does not imply \(a(0)=0\);
intermediate pole corrections must therefore be kept.

\subsection{The complete Weil form and the unresolved identification}

Define the full arithmetic form in these conventions by
\begin{equation}\label{eq:Qdefinition}
Q[F]=\Gamma[F]-W_{1/2}[F]+2\Rea a_F(\ii/2).
\end{equation}
Reflection \(\beta+\ii\gamma\mapsto1-\beta+\ii\gamma\)
and Schwarz reflection of \(a_F\) pair the off-line terms as
conjugates. Thus \eqref{eq:critical-explicit} is equivalently
\begin{equation}\label{eq:full-zero-pairing}
Q[F]=\Zcrit[F]+\Zoff[F]
=\sum_\rho m_\rho
a_F\bigl(\Ima\rho+\ii(\Rea\rho-1/2)\bigr).
\end{equation}
This is the usual complete zero pairing, written in the chosen
Fourier orientation. For sources \eqref{eq:prepared},
\(\widehat F(z)=(z^2+1/4)\widehat h(z)\); hence
\(a_F(\pm\ii/2)=0\) and \(Q[F]=\Gamma[F]-W_{1/2}[F]\).

\begin{corollary}[Remaining arithmetic discrepancy]\label{cor:arithmetic}
For every compact smooth source, with \(Q\) defined by
\eqref{eq:Qdefinition},
\[
B_\sigma[F]-Q[F]
=-R_\sigma[\widehat F]-\Zoff[F]+o_F(1).
\]
In particular this gives \eqref{eq:intro-arithmetic} for prepared
sources.
\end{corollary}
\begin{proof}
Subtract \eqref{eq:full-zero-pairing} from \eqref{eq:deficit}
with \(f=\widehat F\).
\end{proof}

The identity separates two distinct questions. Even if the exact
return trace were shown to vanish, its identified limit would be
\(\Zcrit\). Recovering \(Q\) additionally requires addressing
the off-line pairing. Conversely, under RH that pairing vanishes,
but the results of this paper still do not prove vanishing return
mass for the exact projection.

The all-support quantifier is essential in any use of Weil
positivity: a fixed source family or a fixed support window is
insufficient. The prepared-source version of the criterion is
discussed in \cite[Appendix C]{CCWeil}. No such all-support
positivity conclusion is obtained here.

\section{Scope and further questions}\label{sec:scope}

The estimates establish a finite positive trace family for every
\(\sigma>1/2\), together with an identified critical limit for its
majorant and fixed Abel regularizations. For exact projections they
give an upper limiting measure and locate every possible loss at
one cutoff spectral endpoint. The solvable rational family shows
that the gap between these statements can persist even when each
prelimit cutoff is bounded away from one.

The next question about this particular zeta family is therefore
concrete: determine whether \(R_\sigma[b]\to0\) for every compact
smooth frequency test, or exhibit a test with a nonzero limiting
return defect. Condition~\eqref{eq:uniform-integrability} is an
equivalent endpoint condition once
Theorem~\ref{thm:return} is available. An approximate kernel vector
alone does not settle the question. When nonzero singular values
of \(T_\sigma\) approach zero, a small constraint residual does
not bound distance to the exact kernel.

The phase path studied here is distinct from taking an increasing
finite set of primes at the critical exponent in a semilocal Sonin
construction. In the absolute-convergence region, the causal
transport \eqref{eq:transport} connects the constructions.
No equality of their critical limiting procedures is asserted.
Likewise the two-parameter crossover in
Corollary~\ref{cor:crossover} concerns the rational model; its
scale is not an estimate for zeta.

The theorem package can therefore be assessed independently of
whether this phase family ultimately yields an arithmetic
positivity mechanism. Its content is the source-uniform trace
control, the majorant limit, and the precise distinction between
phase concentration and exact-kernel mass.

\section*{Acknowledgements and use of language models}

This draft was prepared for Edward Baker with substantial assistance
from OpenAI's GPT-6 (Codex). Language-model assistance was used in
the underlying research notes and in consolidating definitions,
developing and checking proofs, organizing the manuscript, and
checking bibliographic links. The exact serving variant and
configured reasoning effort for this drafting session were not
exposed. The internal checks are not independent specialist
refereeing; author verification and approval of the mathematical
claims and final text remain outstanding. No language model is
listed as an author.

\begin{thebibliography}{99}

\bibitem{ConnesTrace}
A.~Connes,
\emph{Trace formula in noncommutative geometry and the zeros of the
Riemann zeta function},
Selecta Math.\ (N.S.) \textbf{5} (1999), 29--106.
\href{https://arxiv.org/abs/math/9811068}{arXiv:math/9811068}.

\bibitem{CCWeil}
A.~Connes and C.~Consani,
\emph{Weil positivity and Trace formula, the archimedean place},
2020 preprint.
\href{https://arxiv.org/abs/2006.13771}{arXiv:2006.13771}.
Theorem~7 and Appendices C--E are used for comparison and conventions.

\bibitem{CCQuasi}
A.~Connes and C.~Consani,
\emph{Quasi-inner functions and local factors},
2020 preprint.
\href{https://arxiv.org/abs/2008.10974}{arXiv:2008.10974}.

\bibitem{CCM}
A.~Connes, C.~Consani, and H.~Moscovici,
\emph{Zeta zeros and prolate wave operators},
2023; version 2, 2024.
\href{https://arxiv.org/abs/2310.18423v2}{arXiv:2310.18423v2}.
Sections 4.6--4.8 concern semilocal Sonin transport.

\bibitem{Burnol}
J.-F.~Burnol,
\emph{Two complete and minimal systems associated with the zeros of
the Riemann zeta function},
2002 preprint.
\href{https://arxiv.org/abs/math/0203120}{arXiv:math/0203120}.

\bibitem{Halmos}
P.~R.~Halmos,
\emph{Two subspaces},
Trans.\ Amer.\ Math.\ Soc.\ \textbf{144} (1969), 381--389.
\href{https://doi.org/10.1090/S0002-9947-1969-0251519-5}
{doi:10.1090/S0002-9947-1969-0251519-5}.

\bibitem{Tao}
T.~Tao,
\emph{254A, Notes 2: Complex-analytic multiplicative number theory},
9 December 2014, Propositions 16 and 19.
\href{https://terrytao.wordpress.com/2014/12/09/254a-notes-2-complex-analytic-multiplicative-number-theory/}
{Author's lecture notes}.

\bibitem{Elkies}
N.~D.~Elkies,
\emph{Math 259: Introduction to Analytic Number Theory},
notes on the Riemann zeta function, pp.~2--4.
\href{https://people.math.harvard.edu/~elkies/M259.02/zeta2.pdf}
{Author's lecture notes}.

\bibitem{DLMF}
F.~W.~J.~Olver et al., eds.,
\emph{NIST Digital Library of Mathematical Functions},
Chapter~25, especially Sections 25.2 and 25.4.
\href{https://dlmf.nist.gov/25.4}{Functional equations of the Riemann zeta function}.

\end{thebibliography}
\end{document}
